% ---
% id: foundations/lem:pinchmono
% kind: lem
% title: "Pinching lowers the trace power"
% summary: >-
%   For α > 1, pinching the α-th root of a positive semidefinite operator and
%   raising the result to the power α never increases the trace, and leaves it
%   unchanged exactly when the operator is block diagonal for the pinching.
%   This bounds the trace-power terms of the MEAT objective.
% prominence: main
% status: implemented
% uses: [def:pinching]
% assumes: []
% refs: []
% provenance: newly-derived
% implements: certiqs_key.meat_core.objective.branch:trace_power_pinched
% lean:
% notes: "The proof uses no result beyond the operator conventions and the pinching of def:pinching, the equality case included."
% ---

\begin{lemma}\label{lem:pinchmono}
Let $\tau\in\Pos{Q}$, let $\cZ$ be a pinching with projectors
$\{Z_j\}_{j=1}^{m}$, and let $\alpha\in(1,\infty)$. Then
\begin{equation}\label{eq:pinchmono}
  \tr\!\left[\left(\cZ\big(\tau^{1/\alpha}\big)\right)^{\alpha}\right]
  \;\le\;\tr[\tau],
\end{equation}
with equality if and only if $\cZ(\tau^{1/\alpha})=\tau^{1/\alpha}$,
equivalently iff $\tau$ is block-diagonal for $\{Z_j\}$, equivalently iff
$\cZ(\tau)=\tau$.
\end{lemma}

\begin{proof}
\emph{Step 1 (pinching as a unitary mixture).} With
$\zeta\coloneqq e^{2\pi i/m}$ and the unitaries
$U_k\coloneqq\sum_{j=1}^{m}\zeta^{jk}Z_j$, $k\in\{0,\dots,m-1\}$,
\eqref{eq:pinch-unitary-mixture} of Definition~\ref{def:pinching} gives
$\cZ(X)=\frac1m\sum_{k=0}^{m-1}U_kXU_k^{\dagger}$ for every $X$.

\emph{Step 2 (strict convexity of $A\mapsto\tr[A^{\alpha}]$ on
$\Pos{Q}$).} Two standard facts, proved inline for self-containment. (a)
\emph{Peierls inequality with equality condition:} for convex
$f:[0,\infty)\to\mathbb{R}$, any ONB $\{\ket{i}\}$ and $A\in\Pos{Q}$ with
spectral decomposition $A=\sum_j\lambda_j\dyad{a_j}$,
\begin{equation*}
  \bra{i}f(A)\ket{i}
  =\sum_j|\braket{i}{a_j}|^2f(\lambda_j)
  \;\ge\;f\Big(\sum_j|\braket{i}{a_j}|^2\lambda_j\Big)
  =f\big(\bra{i}A\ket{i}\big)
\end{equation*}
by Jensen; summing over $i$ gives
$\tr f(A)\ge\sum_if(\bra{i}A\ket{i})$. If $f$ is strictly convex, equality
for index $i$ forces the distribution $|\braket{i}{a_j}|^2$ to be supported
on a single eigenvalue of $A$, i.e.\ $\ket{i}$ lies in one eigenspace of
$A$; equality for all $i$ forces $\{\ket{i}\}$ to be an eigenbasis of $A$.
(b) \emph{Strict convexity of the trace functional:} let
$f(t)=t^{\alpha}$, strictly convex on $[0,\infty)$ for $\alpha>1$. For
$A,B\in\Pos{Q}$ let $C=\tfrac12(A+B)$ with eigenbasis $\{\ket{i}\}$. Then
\begin{equation*}
  \tr f(C)=\sum_if\big(\bra{i}C\ket{i}\big)
  =\sum_if\Big(\tfrac{\bra iA\ket i+\bra iB\ket i}{2}\Big)
  \le\sum_i\tfrac{f(\bra iA\ket i)+f(\bra iB\ket i)}{2}
  \le\tfrac{\tr f(A)+\tr f(B)}{2},
\end{equation*}
using scalar convexity and then Peierls for $A$ and $B$. By continuity,
midpoint convexity gives convexity. If equality holds throughout, then (by
strict scalar convexity) $\bra iA\ket i=\bra iB\ket i$ for all $i$, and (by
the Peierls equality condition) $\{\ket i\}$ is an eigenbasis of both $A$
and $B$; hence $A$ and $B$ are diagonal in a common basis with equal
diagonals, so $A=B$. Thus $A\mapsto\tr[A^{\alpha}]$ is strictly convex.

\emph{Step 3 (conclusion).} Set $X\coloneqq\tau^{1/\alpha}$. By Steps 1--2
and unitary invariance of the trace functional,
\begin{equation*}
  \tr\big[\cZ(X)^{\alpha}\big]
  =\tr\Big[\Big(\tfrac1m\sum_kU_kXU_k^{\dagger}\Big)^{\alpha}\Big]
  \le\frac1m\sum_k\tr\big[(U_kXU_k^{\dagger})^{\alpha}\big]
  =\tr[X^{\alpha}]=\tr[\tau].
\end{equation*}
For the equality case: a strictly convex functional satisfies the
$m$-point Jensen inequality with equality only if all arguments coincide
(if two differed, replacing them by their midpoint would strictly decrease
the right side while fixing the left, a contradiction). Hence equality in
\eqref{eq:pinchmono} forces $U_kXU_k^{\dagger}=U_0XU_0^{\dagger}=X$ for all
$k$, so $X=\frac1m\sum_kU_kXU_k^{\dagger}=\cZ(X)$, i.e.\
$\tau^{1/\alpha}$ is block-diagonal. Since $\tau=(\tau^{1/\alpha})^{\alpha}$
and functional calculus preserves block-diagonality in both directions,
this is equivalent to $\cZ(\tau)=\tau$. Conversely, if
$\cZ(\tau^{1/\alpha})=\tau^{1/\alpha}$ then
\eqref{eq:pinchmono} holds with equality by direct substitution.
\end{proof}
